% Insert as A.81 after the exact one-variable norm tower A.77.
\subsection{The norm class at every integral finite weight}
\begin{statement}{A.81}{Proved}{The norm compatibility class inside the integral finite-weight phantom groups}
Fix a prime $\ell$ and one ordinary pair $(A,B)$ to which
the fixed-level conclusions of A.77 apply.
Keep its compatible root fields $K_n$, with
$[K_n:\Q]=\ell^n$, and put
$Z_n=\ker(\widehat{K_n^\times}_\ell
\xrightarrow{N_{K_n/\Q}}\widehat{\Q^\times}_\ell)$.
The norm compatibility class of A.77 has a canonical
injective image in the full one-variable metabelian
obstruction group.
Its images in every integral finite-weight quotient are
coefficient-$\ell$-power phantom classes.
The proof gives all these maps and their sign.
\end{statement}
\begin{proof}
\emph{The actual kernel in derived coordinates.}
Let $S=X-1$, let $\Lambda=\Z_\ell[[S]]$, and represent
$A,B$ by continuous cocycles
$a,b:G_\Q\to\Z_\ell(1)$.
Set $\theta_g(X)=X^{\chi(g)}$.
Write $V_A$ for the compact module $\Lambda$ with
$g\cdot v=\chi(g)X^{a(g)}\theta_g(v)$.
The exact coefficient sequence governing the actual
one-variable quotient of A.77 is
\[
 0\longrightarrow D_A\xrightarrow{\,S\,}V_A
       \xrightarrow{\epsilon}\Z_\ell(1)\longrightarrow0,
 \qquad
 g\cdot d=\chi(g)X^{a(g)}
       \frac{X^{\chi(g)}-1}{X-1}\theta_g(d)
       \quad(d\in D_A=\Lambda).
 \tag{A.81.1}
\]
Multiplication by $S$ is equivariant by the displayed
action, so the notation distinguishes the derived
generator from the generator $y$ of the permutation
module $V_A$.
The global obstruction is
$o_{\rm pol}(A,B)=\delta[b]\in H^2_{\rm cts}(G_\Q,D_A)$.

Put $N_n=\ell^n$ and
\[
 f_n=X^{N_n}-1,\qquad
 L_n=\frac{X^{N_n}-1}{X-1},\qquad
 V_n=V_A/f_nV_A,\qquad
 J_n=\ker(V_n\xrightarrow{\epsilon}\Z_\ell(1)).
 \tag{A.81.2}
\]
The finite root-torsor map of A.77 is precisely the
coefficient quotient $V_A\to V_n$.
Since $f_n=SL_n$, multiplication by $S$ induces the
actual equivariant isomorphism
\[
                         D_A/L_nD_A\xrightarrow{\ \sim\ }J_n.
 \tag{A.81.3}
\]
Thus the kernel in the rank-one derived coordinate is
$L_nD_A$.
These ideals are stable: the quotient
$\theta_g(L_n)/L_n$ is the ratio of the integral
units
$L_{X^{N_n}}(\chi(g))$ and $L_X(\chi(g))$,
both with constant coefficient $\chi(g)$.

As compact $\Z_\ell$ modules, $J_n$ is free of rank
$N_n-1$, using the monic polynomial $L_n$.
The transition $J_{n+1}\to J_n$ is onto, because
$L_n$ divides $L_{n+1}$.
These maps split continuously as underlying
$\Z_\ell$ modules, without an equivariance requirement.
Moreover
\[
                         D_A\simeq\varprojlim_nJ_n.
 \tag{A.81.4}
\]
Injectivity follows because $L_n$ tends to zero in
the $(\ell,S)$-adic topology.
Surjectivity follows from compactness of the compatible
cosets of the nested ideals $L_n\Lambda$.
The same linear splittings make the inverse systems
of continuous cochains degreewise surjective.

\emph{Norm-one kernels and the precise Milnor injection.}
The integral invariant group
$H^0(G_\Q,\Z_\ell(1))$ is zero.
The number fields $K_n$ likewise have no compatible
roots of unity of arbitrarily high $\ell$-power order,
so $H^0(G_{K_n},\Z_\ell(1))=0$.
By the actual root-torsor Shapiro map of A.77,
$H^1(G_\Q,V_n)=\widehat{K_n^\times}_\ell$,
and augmentation is the norm to $\Q$.
The continuous long exact sequence for
$0\to J_n\to V_n\to\Z_\ell(1)\to0$
therefore gives
\[
 H^1_{\rm cts}(G_\Q,J_n)
       \xrightarrow{\ \sim\ }Z_n .
 \tag{A.81.5}
\]
This isomorphism commutes with the quotient transitions,
which are the field norm maps on the right.
The Milnor sequence from (A.81.4) now gives
\[
 0\longrightarrow\varprojlim\nolimits_n^1 Z_n
   \xrightarrow{\ \iota\ }
   H^2_{\rm cts}(G_\Q,D_A)
   \longrightarrow
   \varprojlim_n H^2_{\rm cts}(G_\Q,J_n)
   \longrightarrow0 .
 \tag{A.81.6}
\]
The algebraic convention is
$\Delta(v)_n=v_n-\pi_n(v_{n+1})$ on additive cochains,
equivalently $z_n/Nz_{n+1}$ on multiplicative classes.

Choose the fixed-level norm preimages $\beta_n$
of A.77.
Choose continuous one-cocycles $h_n$ in $V_n$
representing them with augmentation exactly $b$.
This equality can be imposed on cochains: an
augmentation discrepancy is a coboundary in
$\Z_\ell(1)$; lift its zero-cochain to $V_n$
and subtract that coboundary.
Then $h_n-\pi_nh_{n+1}$ is a $J_n$-valued
one-cocycle representing
$d_n'=\beta_n/N\beta_{n+1}$.
Choose any continuous set lift $w$ of $b$ to $V_A$,
and let $w_n$ be its projections.
The cochains $v_n=h_n-w_n$ lie in $J_n$ and satisfy
$\Delta(v)_n=h_n-\pi_nh_{n+1}$ and
$\delta v_n=-\delta w_n$.
This is the explicit construction of the connecting
homomorphism in (A.81.6), and it fixes the sign:
\[
                  \iota([d_n'])=-\,o_{\rm pol}(A,B).
 \tag{A.81.7}
\]
The right side has zero image at every $J_n$ by
the actual fixed-level norm sections in A.77.
Thus (A.81.7) identifies their exact compatibility
obstruction with the full continuous class.

\emph{Its images in the integral finite-weight groups.}
Put $D_{A,r}=D_A/S^rD_A$, for $r\geq1$.
Its graded modules are
$\Z_\ell(2),\Z_\ell(3),\ldots,\Z_\ell(r+1)$.
The finite-generation proof in A.75 therefore applies
word for word: every $H^1(G_\Q,D_{A,r})$ is finitely
generated over $\Z_\ell$.
Its continuous coefficient maps are $\Z_\ell$ linear,
so compactness kills the corresponding
$S$-adic $\varprojlim^1$ and gives
\[
 H^2_{\rm cts}(G_\Q,D_A)
       \xrightarrow{\ \sim\ }
       \varprojlim_r H^2_{\rm cts}(G_\Q,D_{A,r}).
 \tag{A.81.8}
\]
For $1\leq j\leq\ell^n$ the elementary binomial identity
gives
\[
 v_\ell\binom{\ell^n}{j}=n-v_\ell(j),\qquad
 L_n=\sum_{j=1}^{\ell^n}\binom{\ell^n}{j}S^{j-1}.
 \tag{A.81.9}
\]
Indeed
$\binom{\ell^n}{j}=(\ell^n/j)
\binom{\ell^n-1}{j-1}$, and the last factor is an
$\ell$-adic unit, as is seen by reducing
$(1+S)^{\ell^n-1}$ modulo $\ell$.
For fixed $r,k$, (A.81.9) shows that
$L_n\in(\ell^k,S^r)$ as soon as
$n\geq k+\lfloor\log_\ell r\rfloor$ and
$\ell^n\geq r$.
Hence there is an actual equivariant coefficient map
$J_n\to D_{A,r}/\ell^kD_{A,r}$.
The image of (A.81.7) under this map is zero by
(A.81.6).

To specify the remaining kernels, define
\[
 \mathcal P_{\ell,r}(A)=
 \ker\left(
 H^2_{\rm cts}(G_\Q,D_{A,r})\longrightarrow
 \varprojlim_k
 H^2(G_\Q,D_{A,r}/\ell^kD_{A,r})
 \right).
 \tag{A.81.10}
\]
The transition maps of the finite coefficient
modules have continuous set sections.  These
lift continuous cochains, so their cochain
systems are degreewise surjective and again give
the exact sequence
\[
 0\longrightarrow
 \varprojlim\nolimits_k^1
 H^1(G_\Q,D_{A,r}/\ell^kD_{A,r})
 \longrightarrow H^2_{\rm cts}(G_\Q,D_{A,r})
 \longrightarrow
 \varprojlim_kH^2(G_\Q,D_{A,r}/\ell^kD_{A,r})
 \longrightarrow0 .
 \tag{A.81.11}
\]
Thus $\mathcal P_{\ell,r}(A)$ is exactly the image
of the indicated coefficient-power $\varprojlim^1$.
Every finite-weight image of $\iota([d_n'])$
belongs to $\mathcal P_{\ell,r}(A)$.
Combined with the injection (A.81.6) and
isomorphism (A.81.8), this is an injective
identification of the norm compatibility class
with this compatible family of coefficient-power
phantoms for the same fixed pair.
For $r=1$ its image is zero for an ordinary pair
by A.27 and the actual class-two comparison A.38.
For a pair furnished with a fixed class-$R$ lift
by A.65, the images through $r\leq R-1$ are zero.
No argument here sets all these images to zero
for one fixed ordinary ghost.
\end{proof}

The exact certificate
\path{code/norm_phantom_certificate.py} and its output
verify 1131 binomial coefficient valuations, 554
finite-weight coefficient projection bounds, and nine
derived-kernel and norm-tower divisibility identities.
